

%% =================================================================
%% [분량보존 복원] 부록 E Golden Compare 7시나리오·47점 점별 결과
%% =================================================================
\subsection{7개 시나리오 $\DRTO$ 계산 상세}
\label{app:golden-scenarios}

\subsubsection{계산 프레임워크}
\label{app:golden-framework}

각 시나리오에 대해 고정된 입력 매개변수로 이론적 $\DRTO$를 산출한다.
$\DRTO$ 개별 시그모이드 바운딩 모델은 \m{식~(\ref{eq:app-golden-model})에 나타낸 바와 같다.}

\begin{equation}
  \DRTO = R_0
  + \underbrace{-R_0 \cdot S\!\bigl(\kappa \cdot k_O \cdot O_{\text{raw}}\bigr)}_{E_{O,bnd}}
  + \underbrace{(R_{\max} - R_0) \cdot S\!\bigl(\kappa \cdot k_U \cdot \tfrac{R_0}{R_{\max}-R_0} \cdot U_{\text{raw}}\bigr)}_{E_{U,bnd}}
  \label{eq:app-golden-model}
\end{equation}

이고, 여기서 $S(z) = 2/(1+e^{-z})-1$, $\kappa = 1.2$, $k_U = 1 - k_O$이다.

\subsubsection{시나리오 정의}
\label{app:golden-scenario-def}

7개 산업 시나리오의 매개변수를 \p{[}표~\ref{tab:app-golden-scenarios}\p{]}에 정의한다.

\begin{table}[htbp]
\centering
\caption{Golden Compare 7개 검증 시나리오 매개변수}
\label{tab:app-golden-scenarios}
\small
\begin{tabular}{clccccl}
\toprule
\textbf{\#} & \textbf{시나리오} & \textbf{$R_0$ (h)} & \textbf{$R_{\max}$ (h)}
  & \textbf{$O_{\text{raw}}$} & \textbf{$k_O$} & \textbf{특성} \\
\midrule
1 & \m{기준선(Baseline)}  & 6.0  & 24.0 & 0.50 & 0.50 & 기본 파라미터 \\
2 & 제조업             & 4.0  & 16.0 & 0.60 & 0.55 & 중간 $R_0$, 균형적 관측성 \\
3 & 금융               & 1.0  & 4.0  & 0.77 & 0.45 & 짧은 $R_0$, 높은 관측성 \\
4 & 의료               & 2.0  & 8.0  & 0.68 & 0.50 & 짧은 $R_0$, 높은 관측 필요 \\
5 & IT서비스            & 0.5  & 2.0  & 0.82 & 0.55 & 가장 짧은 $R_0$, 최고 관측성 \\
6 & 물류               & 6.0  & 24.0 & 0.53 & 0.48 & 긴 $R_0$, 중간 관측성 \\
7 & 소매               & 3.0  & 12.0 & 0.57 & 0.52 & 중간 규모, 균형적 \\
\bottomrule
\end{tabular}
\end{table}

\chg{모든 시나리오에서 $R_{\max} = 4R_0$(MTPD 기반 상한)를 적용하며, 시나리오별 $k_O$와
$O_{\text{raw}}$는 산업 특성을 반영한 대표값이다.}

\subsubsection{시나리오 1의 기준선(Baseline) 계산 상세}
\label{app:golden-baseline}

\chg{시나리오 1(기준선)의 매개변수는 $R_0 = 6.0$h, $R_{\max} = 24.0$h, $O_{\text{raw}} = 0.50$,
$k_O = 0.50$이며, 조건별 계산은 다음과 같다.}

\chg{먼저 C0(정적 기준선)에서는 $\DRTO$가 다음과 같이 기준선과 일치한다.}
\begin{equation*}
  \DRTO_{C0} = R_0 = 6.000\;\text{h}, \quad \eta_{C0} = 1.000
\end{equation*}

\chg{C1(관측성만)의 계산은 다음과 같다.}
\begin{align*}
  z_O &= \kappa \cdot k_O \cdot O_{\text{raw}} = 1.2 \times 0.50 \times 0.50 = 0.300 \\
  S(z_O) &= \frac{2}{1+e^{-0.300}} - 1 = 0.1489 \\
  E_{O,bnd} &= -6.0 \times 0.1489 = -0.894\;\text{h} \\
  \DRTO_{C1} &= 6.0 + (-0.894) + 0 = 5.106\;\text{h}, \quad \eta_{C1} = 0.851
\end{align*}

\chg{C2(가우시안 불확실성 추가, $U_{\text{raw}} = 3.0$ 대표값, $k_U = 0.5$)의 계산은 다음과 같다.}
\begin{align*}
  z_U &= \kappa \cdot k_U \cdot \frac{R_0 \cdot U_{\text{raw}}}{R_{\max} - R_0}
       = 1.2 \times 0.5 \times \frac{6.0 \times 3.0}{18.0} = 0.600 \\
  S(z_U) &= \frac{2}{1+e^{-0.600}} - 1 = 0.2913 \\
  E_{U,bnd} &= (24.0 - 6.0) \times 0.2913 = 5.244\;\text{h} \\
  \DRTO_{C2} &= 6.0 + (-0.894) + 5.244 = 10.350\;\text{h}, \quad \eta_{C2} = 1.725
\end{align*}

\chg{C3(파레토 불확실성 추가, $U_{\text{raw}} = 3.0$ 대표값, $k_U = 0.5$)에서는} 파레토의 대표값도
기댓값 $E[U_{\text{raw}}] = 3.0$을 사용하므로 $z_U$ 산출이 C2와 동일하다.
\begin{align*}
  z_U &= \kappa \cdot k_U \cdot \frac{R_0 \cdot U_{\text{raw}}}{R_{\max} - R_0}
       = 1.2 \times 0.5 \times \frac{6.0 \times 3.0}{18.0} = 0.600 \\
  S(z_U) &= \frac{2}{1+e^{-0.600}} - 1 = 0.2913 \\
  E_{U,bnd} &= (24.0 - 6.0) \times 0.2913 = 5.244\;\text{h} \\
  \DRTO_{C3} &= 6.0 + (-0.894) + 5.244 = 10.350\;\text{h}, \quad \eta_{C3} = 1.725
\end{align*}

대표값(평균)에서는 C2와 C3의 $\DRTO$가 일치한다.
두 조건의 차이는 평균이 아닌 \textbf{꼬리 영역}에서 발현된다.
파레토의 P99($U_{\text{raw}} = 21.85$)를 대입하면
$z_U = 1.2 \times 0.5 \times 21.85/3 = 4.370$,
$S(4.370) = \erf{0.9875}{0.975}$,
$E_{U,bnd} = 18.0 \times \erf{0.9875}{0.975} = \erf{17.775}{17.55}$h,
$\DRTO = 6.0 - 0.894 + \erf{17.775}{17.55} = \erf{22.881}{22.657}$h ($\eta = \erf{3.814}{3.776}$)로
C2의 P99($\DRTO = 18.35$h)보다 현저히 높다.

\chg{검증 결과} $0 < \text{C1}(5.106) < \text{C2}(10.350) < \text{C3}_{\text{P99}}(\erf{22.881}{22.657}) < 24.0$\chg{으로 물리적 경계를 충족한다.}

\subsubsection{산업별 C0, C1, C2, C3 계산 요약}
\label{app:golden-industry-c1}

7개 시나리오에 대해 C0(정적 기준선), C1(관측성), C2/C3(불확실성 추가) 조건의
$\DRTO$를 \p{[}표~\ref{tab:app-golden-c1-summary}\p{]}에 정리하였다.
C2와 C3는 등평균 설계($E[U_{\text{raw}}] = 3.0$)이므로
대표값 계산에서는 동일한 결과를 보인다.

\begin{table}[htbp]
\centering
\caption{산업별 $\DRTO$ 계산 \m{결과(대표값 기준)}}
\label{tab:app-golden-c1-summary}
\small
\begin{tabular}{clrrrrrrrr}
\toprule
\textbf{\#} & \textbf{시나리오}
  & \textbf{$\DRTO_{C0}$} & \textbf{$\DRTO_{C1}$} & \textbf{$\eta_{C1}$}
  & \textbf{감소율}
  & \textbf{$\DRTO_{\text{C2=C3}}$} & \textbf{$\eta_{\text{C2=C3}}$}
  & \textbf{$\DRTO_{C3,\text{P99}}$} & \textbf{$\eta_{C3,\text{P99}}$} \\
  & & (h) & (h) & & & (h) & & (h) & \\
\midrule
1 & 기준선    & 6.000 & 5.107 & 0.851 & 14.9\% & 10.350 & 1.725 & 22.657 & 3.776 \\
2 & 제조업    & 4.000 & 3.218 & 0.805 & 19.5\% &  6.382 & 1.595 & 14.757 & 3.689 \\
3 & 금융      & 1.000 & 0.795 & 0.795 & 20.5\% &  1.751 & 1.751 &  3.746 & 3.746 \\
4 & 의료      & 2.000 & 1.598 & 0.799 & 20.1\% &  3.345 & 1.673 &  7.448 & 3.724 \\
5 & IT서비스   & 0.500 & 0.368 & 0.736 & 26.4\% &  0.763 & 1.527 &  1.810 & 3.620 \\
6 & 물류      & 6.000 & 5.091 & 0.849 & 15.1\% & 10.532 & 1.755 & 22.713 & 3.786 \\
7 & 소매      & 3.000 & 2.472 & 0.824 & 17.6\% &  4.995 & 1.665 & 11.205 & 3.735 \\
\bottomrule
\end{tabular}
\end{table}

\chg{표의 C1은 $E_{O,bnd} = -R_0 \cdot S(\kappa \cdot k_O \cdot O_{\text{raw}})$($\kappa = 1.2$)로 계산하며,
C2와 C3는 $U_{\text{raw}} = 3.0$(대표값, 두 분포의 공통 평균)에서 계산하므로 등평균 설계상 대표값에서는
두 조건이 일치한다. $\DRTO_{C3,\text{P99}}$는 파레토 P99($U_{\text{raw}} = 21.85$) 대입값으로, 꼬리 영역에서
C3가 C2보다 현저히 높은 $\DRTO$를 유발함을 보여준다.}


\subsection{47개 검증 포인트 전체 결과}
\label{app:golden-47points}

Golden Compare의 검증 포인트는 4개 카테고리, 총 47개로 구성된다.
\chg{47개의 산출 근거는 다음과 같다. 카테고리 1(순서 제약)은 7개 시나리오에 P99 기준
$C1 \leq C2 \leq C3$ 1건씩을 적용하여 7개, 카테고리 2(물리적 경계)는 7개 시나리오에 4개
조건(C0, C1, C2, C3)을 적용하여 28개, 카테고리 3($z_O$ 동일성)은 6개 시나리오에 C1=C2=C3 확인
1건씩을 적용하여 6개, 카테고리 4(P99 비율 일관성)는 6개 시나리오에 C3/C2 비율 1건씩을 적용하여
6개로 구성된다. 카테고리 3과 4의 시나리오 수가 6개인 것은 시나리오 1(기준선)이 상세 계산에서 별도로
검증되어 시나리오 2--7만을 대상으로 하기 때문이다. 따라서 합계는 47개이며, 모든
포인트에서 PASS를 달성하였다.}

\subsubsection{카테고리 1의 순서 제약(7 포인트)}
\label{app:golden-order}

\p{[}표~\ref{tab:app-golden-order}\p{]}은 P99 기준에서 $\DRTO_{C1} \leq \DRTO_{C2} \leq \DRTO_{C3}$ 순서가 7개 시나리오 모두에서 성립하는지 검증한 결과를 보여준다.

\begin{table}[htbp]
\centering
\caption{순서 제약 검증 \m{결과(P99 기준)}}
\label{tab:app-golden-order}
\small
\begin{tabular}{clcccc}
\toprule
\textbf{\#} & \textbf{시나리오}
  & \textbf{$\DRTO_{C1,\text{P99}}$ (h)} & \textbf{$\DRTO_{C2,\text{P99}}$ (h)}
  & \textbf{$\DRTO_{C3,\text{P99}}$ (h)} & \textbf{판정} \\
\midrule
1 & 기준선    & 5.998 & 8.443  & 13.224 & PASS \\
2 & 제조업    & 3.997 & 5.670  & 8.872  & PASS \\
3 & 금융      & 0.999 & 1.398  & 2.218  & PASS \\
4 & 의료      & 1.998 & 2.830  & 4.436  & PASS \\
5 & IT서비스   & 0.500 & 0.704  & 1.109  & PASS \\
6 & 물류      & 5.998 & 8.443  & 13.224 & PASS \\
7 & 소매      & 2.998 & 4.221  & 6.612  & PASS \\
\midrule
\multicolumn{5}{l}{\textbf{합계: 7/7 PASS}} & \textbf{100\%} \\
\bottomrule
\end{tabular}
\end{table}

\subsubsection{카테고리 2의 물리적 경계(28 포인트)}
\label{app:golden-boundary}

\p{[}표~\ref{tab:app-golden-boundary}\p{]}는 7개 시나리오 $\times$ 4개 조건 $= 28$ 개 $\DRTO$ 값 모두에서 $0 \leq \DRTO \leq R_{\max}$가 성립하는지 검증한 결과를 보여준다.

\begin{table}[htbp]
\centering
\caption{물리적 경계 검증 \m{결과(전체 28 포인트)}}
\label{tab:app-golden-boundary}
\small
\begin{tabular}{clrrrrl}
\toprule
\textbf{\#} & \textbf{시나리오} & \textbf{$R_{\max}$ (h)}
  & \textbf{C0} & \textbf{C1} & \textbf{C2/C3 최대} & \textbf{경계 이내} \\
\midrule
1 & 기준선    & 24.0 & 6.000 & 5.106 & 13.224 & PASS (4건) \\
2 & 제조업    & 16.0 & 4.000 & 3.221 & 8.872  & PASS (4건) \\
3 & 금융      & 4.0  & 1.000 & 0.796 & 2.218  & PASS (4건) \\
4 & 의료      & 8.0  & 2.000 & 1.599 & 4.436  & PASS (4건) \\
5 & IT서비스   & 2.0  & 0.500 & 0.369 & 1.109  & PASS (4건) \\
6 & 물류      & 24.0 & 6.000 & 5.095 & 13.224 & PASS (4건) \\
7 & 소매      & 12.0 & 3.000 & 2.474 & 6.612  & PASS (4건) \\
\midrule
\multicolumn{5}{l}{\textbf{합계: 28/28 PASS}} & & \textbf{100\%} \\
\bottomrule
\end{tabular}
\end{table}

\subsubsection{카테고리 3의 $z_O$ 동일성(6 포인트)}
\label{app:golden-zo}

동일 시나리오 내에서 C1, C2, C3의 관측성 입력 $z_O = \kappa \cdot k_O \cdot O_{\text{raw}}$가
조건에 무관하게 동일한 값을 갖는지 검증한다.
$\DRTO$ 모델에서 관측성 채널($E_{O,bnd}$)과 불확실성 채널($E_{U,bnd}$)은
독립적으로 작동하도록 설계되었으므로,
C2/C3에서 불확실성 분포를 다르게 적용하더라도
관측성 쪽 계산($z_O$)은 C1과 완전히 동일해야 한다.
이 조건이 성립해야 ``불확실성 조건을 추가해도 관측성 채널이
오염되지 않는다''는 \textbf{채널 독립성}이 보장된다. \p{[}표~\ref{tab:app-golden-zo}\p{]}는 $z_O$ 동일성 검증 결과를 보여준다.

\begin{table}[htbp]
\centering
\caption{$z_O$ 동일성 검증 결과}
\label{tab:app-golden-zo}
\small
\begin{tabular}{clcccc}
\toprule
\textbf{\#} & \textbf{시나리오}
  & \textbf{$z_O^{(C1)}$} & \textbf{$z_O^{(C2)}$} & \textbf{$z_O^{(C3)}$} & \textbf{판정} \\
\midrule
2 & 제조업    & 0.396 & 0.396 & 0.396 & PASS \\
3 & 금융      & 0.416 & 0.416 & 0.416 & PASS \\
4 & 의료      & 0.408 & 0.408 & 0.408 & PASS \\
5 & IT서비스   & 0.541 & 0.541 & 0.541 & PASS \\
6 & 물류      & 0.306 & 0.306 & 0.306 & PASS \\
7 & 소매      & 0.356 & 0.356 & 0.356 & PASS \\
\midrule
\multicolumn{5}{l}{\textbf{합계: 6/6 PASS}} & \textbf{100\%} \\
\bottomrule
\end{tabular}
\end{table}

\chg{시나리오 1(기준선)은 대표 시나리오로서 별도로 검증하며, $z_O = \kappa \cdot k_O \cdot O_{\text{raw}}$는
조건(C1, C2, C3)에 무관하게 동일하다.}

\subsubsection{카테고리 4의 P99 비율 일관성(6 포인트)}
\label{app:golden-ratio}

$\DRTO_{C3,\text{P99}} / \DRTO_{C2,\text{P99}}$ 비율이
시나리오에 무관하게 안정적인 값($\approx 1.57$)을 유지하는지 검증한다.
원시 입력 $U_{\text{raw}}$의 P99 비율은 파레토/가우시안 $\approx 3.95$이나,
시그모이드 바운딩을 거치면서 비선형 압축이 적용되어
$\DRTO$ 수준에서의 비율은 $\approx 1.57$로 수렴한다(\p{[}표~\ref{tab:app-golden-ratio}\p{]}).

\begin{table}[htbp]
\centering
\caption{P99 \m{비율($\DRTO_{C3,\text{P99}} / \DRTO_{C2,\text{P99}}$)} 일관성 검증}
\label{tab:app-golden-ratio}
\small
\begin{tabular}{clcccl}
\toprule
\textbf{\#} & \textbf{시나리오}
  & \textbf{$\DRTO_{C2,\text{P99}}$} & \textbf{$\DRTO_{C3,\text{P99}}$}
  & \textbf{비율} & \textbf{판정} \\
\midrule
2 & 제조업    & 5.670  & 8.872  & 1.56 & PASS \\
3 & 금융      & 1.398  & 2.218  & 1.59 & PASS \\
4 & 의료      & 2.830  & 4.436  & 1.57 & PASS \\
5 & IT서비스   & 0.704  & 1.109  & 1.58 & PASS \\
6 & 물류      & 8.443  & 13.224 & 1.57 & PASS \\
7 & 소매      & 4.221  & 6.612  & 1.57 & PASS \\
\midrule
\multicolumn{4}{l}{\textbf{합계: 6/6 PASS}} & \textbf{평균 1.57} & \textbf{100\%} \\
\bottomrule
\end{tabular}
\end{table}

\chg{P99 비율의 평균은 $\approx 1.57$로 시나리오 간 높은 일관성을 보이며, 이 비율은 파레토 분포의 P99와
가우시안 P99의 구조적 차이를 반영하여 $R_0$와 $R_{\max}$의 \m{절댓값}에 무관하게 안정적이다.}

\subsubsection{검증 포인트 종합}
\label{app:golden-summary}

4개 카테고리 47개 검증 포인트의 종합 결과를 \p{[}표~\ref{tab:app-golden-total}\p{]}에 요약한다.

\begin{table}[htbp]
\centering
\caption{Golden Compare 47개 검증 포인트 종합 결과}
\label{tab:app-golden-total}
\begin{tabular}{clcl}
\toprule
\textbf{카테고리} & \textbf{검증 내용} & \textbf{포인트 수} & \textbf{결과} \\
\midrule
순서 제약       & $\DRTO_{C1} \leq \DRTO_{C2} \leq \DRTO_{C3}$ (P99 기준)      & 7  & 100\% PASS \\
물리적 경계     & $0 \leq \DRTO \leq R_{\max}$                                  & 28 & 100\% PASS \\
$z_O$ 동일성   & $z_O^{(C1)} = z_O^{(C2)} = z_O^{(C3)}$ (조건 독립)            & 6  & 100\% PASS \\
P99 비율 일관성 & $\DRTO_{C3,\text{P99}} / \DRTO_{C2,\text{P99}} \approx 1.57$           & 6  & 100\% PASS \\
\midrule
\textbf{합계}  &                                                               & \textbf{47} & \textbf{100\% PASS} \\
\bottomrule
\end{tabular}
\end{table}

\chg{47개 검증 포인트 모두에서 PASS를 달성하여, $\DRTO$ 개별 시그모이드 바운딩 모델의 구현 정확성과
SSOT(Single Source of Truth) 일관성이 검증되었다.}

\chg{Golden Compare 검증은 네 가지 결론을 뒷받침한다. 첫째, 구현 정확성으로서 7개 시나리오와 4개
조건을 곱한 28개 $\DRTO$ 값 모두에서 이론적 계산치와 시뮬레이션 출력이 일치한다. 둘째, 물리적
타당성으로서 모든 $\DRTO$가 $(0,\; R_{\max})$ 범위 내에 있어 명제 2.1(물리적 경계
보장)을 실증적으로 확인한다. 셋째, 구조적 일관성으로서 $z_O$ 동일성과 P99 비율 일관성을 통해 관측성과
불확실성 채널의 독립 작동이 검증된다. 넷째, 산업 범용성으로서 기준선과 6개 산업에서 동일한 구조적
패턴이 재현되어 $\DRTO$ 프레임워크의 산업 범용 적용성을 확인한다.}

